\section{Timeline and Resources Required}

\subsection{Timeline}

The project timeline spans 26 weeks, organized into 11 major milestones as illustrated in Figure \ref{fig:timeline}. Each milestone encompasses specific tasks designed to ensure systematic progress toward the development and evaluation of the proposed blockchain, deep reinforcement learning, and LLM-based defense framework for SDVN timing attacks.

\subsubsection*{Milestone 1: Project Topic Selection and Initial Planning (Weeks 1-2)}
\begin{itemize}
    \item Review project requirements and scope
    \item Conduct preliminary research on SDVN security challenges
    \item Finalize project topic and obtain supervisor approval
    \item Establish team roles and communication protocols
\end{itemize}

\subsubsection*{Milestone 2: Background Study and Literature Review (Weeks 2-4)}
\begin{itemize}
    \item Study fundamental concepts of SDN and SDVN architectures
    \item Review existing research on control plane timing attacks
    \item Investigate blockchain applications in network security
    \item Explore deep reinforcement learning and LLM technologies
    \item Conduct comprehensive literature survey on related defense mechanisms
\end{itemize}

\subsubsection*{Milestone 3: Project Proposal Submission (Weeks 4-5)}
\begin{itemize}
    \item Finalize problem statement and research objectives
    \item Complete methodology design and timeline planning
    \item Prepare and submit comprehensive project proposal
    \item Incorporate supervisor feedback and revisions
\end{itemize}

\subsubsection*{Milestone 4: Attack Pattern Analysis and Characterization (Weeks 5-7)}
\begin{itemize}
    \item Characterize delay injection attack patterns
    \item Analyze race condition exploitation scenarios
    \item Study timing side-channel attack mechanisms
    \item Investigate synchronization exploit vulnerabilities
    \item Document attack models and threat scenarios
\end{itemize}

\subsubsection*{Milestone 5: Blockchain Component Development (Weeks 7-10)}
\begin{itemize}
    \item Design blockchain architecture for timing event logging
    \item Implement Hyperledger Fabric network setup
    \item Develop smart contracts for timestamp verification
    \item Implement consensus mechanisms for timing integrity
    \item Test blockchain component independently
\end{itemize}

\subsubsection*{Milestone 6: Deep Reinforcement Learning Framework (Weeks 10-13)}
\begin{itemize}
    \item Design DRL agent architecture and state-action space
    \item Implement reward function for timing policy optimization
    \item Develop training pipeline using PyTorch
    \item Train DRL model with simulated attack scenarios
    \item Evaluate and fine-tune model performance
\end{itemize}

\subsubsection*{Milestone 7: Large Language Model Integration (Weeks 13-15)}
\begin{itemize}
    \item Integrate Hugging Face Transformers framework
    \item Develop log analysis and interpretation modules
    \item Implement contextual threat intelligence system
    \item Design alert generation and recommendation engine
    \item Test LLM contextual understanding capabilities
\end{itemize}

\subsubsection*{Milestone 8: System Integration and Hybrid Framework Development (Weeks 15-17)}
\begin{itemize}
    \item Integrate blockchain, DRL, and LLM components
    \item Implement time synchronization protocols (PTP/NTP)
    \item Develop rate limiting and buffering mechanisms
    \item Create unified defense framework interface
    \item Conduct integration testing and debugging
\end{itemize}

\subsubsection*{Milestone 9: NS3 Simulation Implementation (Weeks 17-20)}
\begin{itemize}
    \item Set up NS3 environment and SDVN modules
    \item Implement vehicular mobility models
    \item Configure attack scenarios and network topologies
    \item Deploy integrated defense framework in NS3
    \item Validate simulation accuracy and reliability
\end{itemize}

\subsubsection*{Milestone 10: Performance Evaluation and Analysis (Weeks 20-24)}
\begin{itemize}
    \item Conduct extensive simulation experiments
    \item Measure attack detection rate and false positive rate
    \item Evaluate system latency and throughput metrics
    \item Compare results with existing defense mechanisms
    \item Perform statistical analysis and result validation
    \item Identify areas for optimization and improvement
\end{itemize}

\subsubsection*{Milestone 11: Documentation and Publication (Weeks 24-26)}
\begin{itemize}
    \item Prepare comprehensive final project report
    \item Create presentation materials for demonstration
    \item Write research paper for journal submission
    \item Prepare for final presentation and viva voce
    \item Submit all deliverables and documentation
\end{itemize}

\begin{figure}[htbp]
    \centering
    \includegraphics[width=0.95\textwidth,angle=90]{project_timeline.png}
    \caption{Project Timeline - Gantt chart showing the 26-week schedule with 11 major milestones and their respective durations.}
    \label{fig:timeline}
\end{figure}

\subsection{Resources Required}

Successful implementation of the proposed defense framework requires specific hardware, software, and computational resources as outlined below.

\subsubsection{Hardware Resources}

\begin{table}[htbp]
\centering
\begin{tabular}{|l|l|}
\hline
\textbf{Component} & \textbf{Specification} \\
\hline
Processor & Intel Core i7 / AMD Ryzen 7 (or equivalent) \\
\hline
RAM & 16GB minimum (32GB recommended) \\
\hline
Storage & 200GB free space for simulations and datasets \\
\hline
GPU & NVIDIA GPU with CUDA support (for DRL training) \\
\hline
Network & Stable internet connection for cloud resources \\
\hline
\end{tabular}
\caption{Hardware Requirements}
\label{tab:hardware}
\end{table}

\subsubsection{Software Resources}

\begin{table}[htbp]
\centering
\begin{tabular}{|l|l|l|}
\hline
\textbf{Category} & \textbf{Software/Tool} & \textbf{Cost} \\
\hline
Network Simulation & NS3 (Network Simulator 3) & Free (Open Source) \\
\hline
Programming Language & Python 3.8+ & Free (Open Source) \\
\hline
& C++ Compiler (GCC/Clang) & Free (Open Source) \\
\hline
Blockchain Platform & Hyperledger Fabric 2.x & Free (Open Source) \\
\hline
& Docker \& Docker Compose & Free (Open Source) \\
\hline
Deep Learning & PyTorch Framework & Free (Open Source) \\
\hline
& TensorFlow (optional) & Free (Open Source) \\
\hline
& CUDA Toolkit & Free \\
\hline
LLM Framework & Hugging Face Transformers & Free (Open Source) \\
\hline
Optimization & GurobiPy Academic License & Free (Academic) \\
\hline
Data Processing & NumPy, Pandas & Free (Open Source) \\
\hline
Visualization & Matplotlib, Seaborn & Free (Open Source) \\
\hline
Development Tools & Git, Jupyter Notebook & Free (Open Source) \\
\hline
Documentation & LaTeX & Free (Open Source) \\
\hline
\end{tabular}
\caption{Software and Development Tools}
\label{tab:software}
\end{table}

\subsubsection{Cloud and Computational Resources (Optional)}

For intensive simulations and large-scale experiments, cloud computing resources may be utilized:

\begin{table}[htbp]
\centering
\begin{tabular}{|l|l|}
\hline
\textbf{Resource} & \textbf{Estimated Cost} \\
\hline
Cloud Computing (AWS/GCP/Azure) & \$50-100 for 100-200 compute hours \\
\hline
Cloud Storage & \$10-20 per month \\
\hline
GPU Instances (if needed) & \$100-200 for intensive training \\
\hline
\textbf{Total Estimated Cloud Cost} & \textbf{\$160-320 (project duration)} \\
\hline
\end{tabular}
\caption{Estimated Cloud Computing Costs (Optional)}
\label{tab:cloud}
\end{table}

\subsubsection{Reference and Learning Resources}

\begin{itemize}
    \item \textbf{Academic Databases:} Access to IEEE Xplore, ACM Digital Library, ScienceDirect, and Springer for research papers
    \item \textbf{Documentation:} Official documentation for NS3, Hyperledger Fabric, PyTorch, and Hugging Face
    \item \textbf{Online Courses:} Coursera, edX, and YouTube tutorials for specialized topics
    \item \textbf{Technical Support:} University computing resources and supervisor consultations
\end{itemize}

\subsubsection{Cost Summary}

\begin{table}[htbp]
\centering
\begin{tabular}{|l|r|}
\hline
\textbf{Resource Category} & \textbf{Estimated Cost (USD)} \\
\hline
Hardware (personal computers) & \$0 (using existing equipment) \\
\hline
Open Source Software & \$0 \\
\hline
Academic Software Licenses & \$0 (free academic licenses) \\
\hline
Cloud Computing (optional) & \$160-320 \\
\hline
Research Paper Access & \$0 (university subscriptions) \\
\hline
Miscellaneous & \$50 \\
\hline
\textbf{Total Estimated Budget} & \textbf{\$210-370} \\
\hline
\end{tabular}
\caption{Overall Project Budget Estimate}
\label{tab:budget}
\end{table}

The project is designed to be cost-effective by leveraging open-source tools and academic licenses. The primary costs involve optional cloud computing resources for large-scale simulations and miscellaneous expenses for documentation and presentation materials.